\section{Natural transformations}\label{sec:natural_transformations}

\paragraph{Bifunctors}

\begin{definition}\label{def:product_category}\mimprovised
  We define the \term{product category} \( \prod_{k \in \mscrK} \cat{C}_k \) of the \hyperref[def:indexed_family]{indexed family} \( \seq{ \cat{C}_k }_{k \in \mscrK} \) of categories as follows\fnote{Equivalently, we may consider the \hyperref[def:categorical_diagram]{diagram} \( D: \cat{K} \to \cat{Cat} \), where \( \cat{K} \) is the \hyperref[def:discrete_category]{discrete category} over \( \mscrK \).}:
  \begin{thmenum}[series=def:product_category]
    \thmitem{def:product_category/objects} The \hyperref[def:category/objects]{set of objects} is the \hyperref[def:cartesian_product]{Cartesian product}
    \begin{equation*}
      \obj\parens*{ \prod_{k \in \mscrK} \cat{C}_k } \coloneqq \prod_{k \in \mscrK} \obj(\cat{C}_k).
    \end{equation*}

    \thmitem{def:product_category/morphisms} The \hyperref[def:category/morphisms]{set of morphisms} is similarly the Cartesian product
    \begin{equation*}
      \mor\parens*{ \prod_{k \in \mscrK} \cat{C}_k } \coloneqq \prod_{k \in \mscrK} \mor(\cat{C}_k).
    \end{equation*}

    For the tuple \( \seq{ f_k }_{k \in \mscrK} \) of morphisms, where \( f_k: A_k \to B_k \), the domain and codomain are the tuples \( \seq{ A_k }_{k \in \mscrK} \) and \( \seq{ B_k }_{k \in \mscrK} \), respectively.

    \thmitem{def:product_category/composition} The \hyperref[def:category/composition]{composition operation} of two tuples of morphisms is their elementwise composition.

    \thmitem{def:product_category/identity} The \hyperref[def:category/identity]{identity morphism} on \( \seq{ A_k }_{k \in \mscrK} \) is the corresponding tuple of identity morphisms.
  \end{thmenum}

  We will associate with each product category the following:
  \begin{thmenum}[resume=def:product_category]
    \thmitem{def:product_category/projection} For every index \( m \) in \( \mscrK \), we combine the projection functions of objects and morphisms of \( \cat{C}_k \) into a functor
    \begin{equation*}
      \pi_m: \prod_{k \in \mscrK} \cat{C}_k \to \cat{C}_m.
    \end{equation*}
  \end{thmenum}
\end{definition}
\begin{comments}
  \item Product categories are \hyperref[def:categorical_product]{products} in the \hyperref[def:category_of_small_categories]{category of small categories} --- see \fullref{thm:product_category_universal_property}.

  \item General product categories are rarely constructed explicitly. Definitions for binary product categories (i.e. \( \cat{C} \times \cat{D} \)) can be found in
  \cite[36]{MacLane1998CategoryTheory},
  \cite[def. 2.5.1]{Schubert1972Categories} or
  \cite[def. 1.3.49]{Perrone2024StartingCategoryTheory}.
\end{comments}

\begin{concept}\label{con:functor_argument_variance}\mimprovised
  If the domain of a \hyperref[def:functor]{functor} \( F: \cat{C}_1 \times \cdots \times \cat{C}_n \to \cat{D} \) is a \hyperref[def:product_category]{product category}, we can treat it as a function of several arguments, as per \cref{con:function_arguments}.

  Furthermore, \( F \) can be covariant or \hyperref[def:contravariant_functor]{contravariant} in the \( k \)-th argument, depending on whether \( \cat{C}_k \) is \emph{intended} to be treated as an \hyperref[def:opposite_category]{opposite category}.

  Like the number of arguments, the variance of the arguments is unambiguous in practice. This is easily seen for the \( \hom \)-functors, defined in \cref{def:hom_functor}, which gives rise to \fullref{sec:yoneda_embeddings} and \fullref{sec:adjoint_functors}. Other examples can be found in \cref{ex:def:natural_transformation} for example.
\end{concept}

\begin{definition}\label{def:bifunctor}\mcite[def. 2.5.2]{Schubert1972Categories}
  We call a \hyperref[def:functor]{functor} \( F: \cat{C} \times \cat{D} \to \cat{E} \) from any \hyperref[def:product_category]{product category} \( \cat{C} \times \cat{D} \) a \term{bifunctor}.
\end{definition}

\begin{definition}\label{def:hom_functor}\mcite[37]{Jacobson1989BasicAlgebraII}
  In a \hyperref[def:small_category]{locally small} \hyperref[def:category]{category} \( \cat{C} \), parametrizing the \( \hom \)-sets gives us the \hyperref[def:bifunctor]{bifunctor}
  \begin{equation*}
    \begin{aligned}
      &\hom: \cat{C}^\oppos \times \cat{C} \to \cat{Set}, \\
      &\hom(A, Y) \coloneqq \set{ f \in \mor(\cat{C}) \given \dom(f) = A \T{and} \cod(f) = Y }, \\
      &\hom(f: B \to A, g: X \to Y) \coloneqq (\underbrace{s}_{\mathclap{A \to X}} \mapsto \underbrace{g \bincirc s \bincirc f}_{B \to Y}).
    \end{aligned}
  \end{equation*}

  It acts as expected on objects, i.e. \( \hom(A, Y) \) is the set of morphisms from \( A \) to \( Y \). Its action on morphisms can be visualized as follows:
  \begin{equation*}
    \begin{aligned}
      \includegraphics[page=1]{output/def__hom_functor}
    \end{aligned}
  \end{equation*}

  By fixing the first pair of arguments to some object \( A \) and its identity morphism \( \id_A \), we obtain the covariant functor
  \begin{equation*}
    \hom(A, \anon*): \cat{C} \to \cat{Set},
  \end{equation*}
  while fixing the second pair of arguments gives us the contravariant functor
  \begin{equation*}
    \hom(\anon*, Y): \cat{C}^\oppos \to \cat{Set}.
  \end{equation*}
\end{definition}
\begin{comments}
  \item \hyperref[con:currying]{Uncurrying} the \( \hom \)-bifunctor leads to the Yoneda embeddings we describe in \cref{def:yoneda_embeddings}.
\end{comments}
\begin{defproof}
  For the bifunctor, \ref{def:functor/CF2} can be seen to hold by inspecting the diagram:
  \begin{equation*}
    \begin{aligned}
      \includegraphics[page=2]{output/def__hom_functor}
    \end{aligned}
  \end{equation*}

  The other functor condition \ref{def:functor/CF1} is straightforward to prove.
\end{defproof}

\begin{example}\label{ex:def:hom_functor}
  We list examples related to the \hyperref[def:hom_functor]{\( \hom \)-bifunctor}:
  \begin{thmenum}
    \thmitem{ex:def:hom_functor/vect} Consider the category \hyperref[def:vector_space/category]{\( \cat[\BbbK]{Vect} \)} of small vector spaces over \( \BbbK \).

    For fixed spaces \( \BbbK^m \) and \( \BbbK^n \), the set \( \hom(\BbbK^n, \BbbK^m) \) consists of the corresponding linear functions or, by \cref{rem:linear_operators_and_matrices}, the \( m \times n \) \hyperref[def:array/matrix]{matrices} over \( \BbbK \).

    Then, the \( \hom \)-bifunctor acts on matrices as follows:
    \begin{equation*}
      \hom(\underbrace{A}_{a \times b}, \underbrace{B}_{d \times c})
      =
      (\underbrace{C}_{c \times a} \mapsto \underbrace{BCA}_{d \times b}).
    \end{equation*}

    In particular, when \( a = c \) and \( b = d = 1 \), \( \hom(A, B) \) consists of \hyperref[def:bilinear_form]{bilinear forms}.
  \end{thmenum}
\end{example}

\begin{remark}\label{rem:opposite_morphism_argument_notation}
  As an abstract entity, the \hyperref[def:opposite_category/morphisms]{opposite morphism} \( f^*: B \to A \) of \( f: A \to B \) may be useful when applying \fullref{thm:category_duality}. In most categories of interest, however, the domain and codomain play an asymmetric role. For example, in any \hyperref[def:concrete_category]{concrete category}, the morphism \( f: A \to B \) can be regarded as a function, while its opposite \( f^*: B \to A \) is an abstract entity with no inherent meaning.

  When used as an argument, an opposite morphism usually indicates that the argument is \hyperref[con:functor_argument_variance]{contravariant}. For example, we define the \( \hom \)-bifunctor in \cref{def:hom_functor} as
  \begin{equation*}
    \begin{aligned}
      &\hom: \cat{C}^\oppos \times \cat{C} \to \cat{Set}, \\
      &\hom(A, Y) \coloneqq \set{ f \in \mor(\cat{C}) \given \dom(f) = A \T{and} \cod(f) = Y }, \\
      &\hom(f: B \to A, g: X \to Y) \coloneqq (\underbrace{s}_{\mathclap{A \to X}} \mapsto \underbrace{g \bincirc s \bincirc f}_{B \to Y}).
    \end{aligned}
  \end{equation*}

  Instead of treating the first argument as the abstract entity \( f^*: A \to B \), we simply swap the usual alphabetical order of its arguments and write \( f: B \to A \). This highlights that \( \hom(\anon*, \anon*) \) is contravariant in its first argument, but avoids formal dualities.

  This convention is useful when dealing with \fullref{sec:yoneda_embeddings}.
\end{remark}

\paragraph{Natural transformations}

\begin{definition}\label{def:natural_transformation}\mcite[def. 7.6; def. 7.10]{Awodey2010CategoryTheory}
  Let \( F \) and \( G \) be parallel \hyperref[def:functor]{functors} from \( \cat{C} \) to \( \cat{D} \).

  A \term[ru=естественное преобразование (\cite[def. 3.1]{ЦаленкоШульгейфер1974ОсновыТеорииКатегорий})]{natural transformation} \( \alpha: F \Rightarrow G \) is an \hyperref[def:indexed_family]{indexed family} \( \seq{ \alpha_A }_{A \in \obj(\cat{C})} \) of morphisms in \( \cat{D} \), where \( \alpha_A: F(A) \to G(A) \), and such that, for every morphism \( f: A \to B \) in \( \cat{C} \), the following diagram, called the \term[en=naturality condition (\cite[def. 1.4.1]{Perrone2024StartingCategoryTheory})]{naturality condition}, \hyperref[def:categorical_diagram/commutative]{commutes}:
  \begin{equation}\label{eq:def:natural_transformation/diagram}
    \begin{aligned}
      \includegraphics[page=1]{output/def__natural_transformation}
    \end{aligned}
  \end{equation}

  The morphisms \( \alpha_A \) are called the \term{components} of \( \alpha \). We denote natural transformations by \( \alpha: F \Rightarrow G \) and, when used in diagrams, by
  \begin{equation}\label{eq:def:natural_transformation/notation}
    \begin{aligned}
      \includegraphics[page=2]{output/def__natural_transformation}
    \end{aligned}
  \end{equation}
\end{definition}

\begin{definition}\label{def:natural_isomorphism}\mcite[def. 7.10]{Awodey2010CategoryTheory}
  We say that a \hyperref[def:natural_transformation]{natural transformation} is a \term[ru=естественный изоморфизм (\cite[def. 3.1]{ЦаленкоШульгейфер1974ОсновыТеорииКатегорий})]{natural isomorphism} if all its components are \hyperref[def:morphism_invertibility/isomorphism]{isomorphisms}.
\end{definition}
\begin{comments}
  \item This definition is justified in \cref{thm:natural_isomorphism}, where we show that the isomorphisms in a \hyperref[def:functor_category]{functor category} are the natural transformations whose components are isomorphisms.

  This does not hold more generally for one-sided inverses or cancellable morphisms.
\end{comments}

\begin{example}\label{ex:def:natural_transformation}
  We list examples of \hyperref[def:natural_transformation]{natural transformations}:
  \begin{thmenum}
    \thmitem{ex:def:natural_transformation/cartesian_product} We largely ignore the distinction between the  \hyperref[def:cartesian_product]{Cartesian products} \( (A \times B) \times C \) and \( A \times (B \times C) \). A precise sense in which this distinction really is irrelevant is the existence of a natural isomorphism.

    We can consider them as functors from the \hyperref[def:product_category]{product category} \( \cat{Set}^3 \) to \( \cat{Set} \), namely
    \begin{align*}
      &L: \cat{Set}^3 \to \cat{Set}, \\
      &L(A, B, C) \coloneqq (A \times B) \times C, \\
      &L(f: A \to X, g: B \to Y, h: C \to Z) \coloneqq \smash{\overbrace{\parens[\big]{ \parens[\big]{ (a, b), c } \mapsto \parens[\big]{ (f(a), g(b)), h(c) } }}^{(A \times B) \times C \to (X \times Y) \times Z}},
    \intertext{and}
      &R: \cat{Set}^3 \to \cat{Set}, \\
      &R(A, B, C) \coloneqq A \times (B \times C), \\
      &R(f: A \to X, g: B \to Y, h: C \to Z) \coloneqq \smash{\overbrace{\parens[\big]{ \parens[\big]{ a, (b, c) } \mapsto \parens[\big]{ f(a), (g(b), h(c)) } }}^{A \times (B \times C) \to X \times (Y \times Z)}}.
    \end{align*}

    We can define the natural transformation
    \begin{equation*}
      \begin{aligned}
        &\alpha: L \Rightarrow R, \\
        &\alpha_{(A, B, C)}\parens[\big]{ (a, b), c } \coloneqq \parens[\big]{ a, (b, c) }.
      \end{aligned}
    \end{equation*}

    For every object \( (A, B, C) \) in \( \cat{Set}^3 \), the function \( \alpha_{(A, B, C)} \) translates a tuple from \( (A \times B) \times C \) to \( A \times (B \times C) \).

    The naturality condition is simple to verify. For any triple of functions \( f: A \to X \), \( g: B \to Y \) and \( h: C \to Z \), we want to show that
    \begin{equation*}
      \begin{aligned}
        \includegraphics[page=1]{output/ex__def__natural_transformation}
      \end{aligned}
    \end{equation*}

    Fortunately, this is trivial. Going through the top-right corner gives us
    \begin{equation*}
      \alpha_{(X, Y, Z)}\parens[\big]{ L(f, g, h)\parens[\big]{ (a, b), c } }
      =
      \alpha_{(X, Y, Z)}\parens[\big]{ (f(a), g(b)), h(c) }
      =
      \parens[\big]{ f(a), (g(b), h(c)) },
    \end{equation*}
    while going through the bottom-left corner gives us
    \begin{equation*}
      R(f, g, h)\parens[\big]{ \alpha_{(A, B, C)}\parens[\big]{ (a, b), c } }
      =
      R(f, g, h)\parens[\big]{ a, (b, c) }
      =
      \parens[\big]{ f(a), (g(b), h(c)) }.
    \end{equation*}

    \thmitem{ex:def:natural_transformation/global_elements} For a fixed (small) \hyperref[def:singleton_set]{singleton set}, say \( 1 = \set{ 0 } = \set{ \varnothing } \), the identity functor in \( \cat{Set} \) is naturally isomorphic to the functor \( \hom(1, \anon*) \).

    Indeed, we can define the (obviously bijective) natural transformation
    \begin{equation*}
      \alpha_A(a) \coloneqq (0 \mapsto a).
    \end{equation*}

    Note that \( \alpha_A: A \to \hom(1, A) \) is defined in the same way for every \( A \), but formally these are different functions.\fnote{Although it is possible to define \( \alpha \) over the \hyperref[def:smallest_grothendieck_universe]{smallest Grothendieck universe} \( \mscrU_0 \)}

    We must verify that the following diagram commutes:
    \begin{equation*}
      \begin{aligned}
        \includegraphics[page=2]{output/ex__def__natural_transformation}
      \end{aligned}
    \end{equation*}

    Indeed,
    \begin{equation*}
      \alpha_B(f(a))
      =
      \parens[\big]{ 0 \mapsto f(a) }
      =
      f \bincirc \alpha_A(a).
    \end{equation*}

    This holds more generally in many other categories with a \hyperref[def:universal_objects/terminal]{terminal object}, such as \( \hyperref[def:group/category]{\cat{Grp}} \).

    The \hyperref[def:directed_multigraph/category]{category of small directed multigraphs} is more subtle. As discussed in \cref{ex:def:universal_objects/directed_multigraphs}, any graph with one vertices and a loop is terminal. A homomorphism from such a graph is a vertex with a loop.

    As a special case that somewhat reconciles the two cases above, in \( \hyperref[def:category_of_small_categories]{\cat{Cat}} \), a constant functor from the \hyperref[def:singleton_category]{terminal category} chooses a vertex and, due to \ref{def:functor/CF1}, its identity morphism. The diagonal functor, defined in \cref{def:diagonal_functor}, sends an object to its corresponding constant functor.

    All this is generalized in \cref{def:global_element}.

    \thmitem{ex:def:natural_transformation/boolean} Similarly, the \hyperref[def:subset_characteristic_function]{subset characteristic functions} \( \chi_A \) for small sets can be grouped into a natural isomorphism between the power set functor \( \pow \) in \( \cat{Set} \) (see \cref{ex:def:functor/power_set}) and the functor \( \hom(2, \anon*) \).

    A generalization of this isomorphism formalizes the connection between \( n \)-ary \hyperref[def:relation]{relations} \hyperref[def:boolean_function]{Boolean-valued functions}, which we discuss in \cref{rem:boolean_valued_functions_and_relations}. In this case the functors become
    \small
    \begin{align*}
      &P: \cat{Set}^n \to \cat{Set}, \\
      &P(A_1, \ldots, A_n) \coloneqq \pow(A_1 \times \cdots \times A_n), \\
      &P(f_1: A_1 \to B_1, \ldots, f_n: A_n \to B_n) \coloneqq \parens[\big]{ S \mapsto \set[\big]{ \parens[\big]{ f_1(a_1), \ldots, f_n(a_n) } \given (a_1, \ldots, a_n) \in S } }
    \intertext{and}
      &B: \cat{Set}^n \to \cat{Set}, \\
      &B(A_1, \ldots, A_n) \coloneqq \hom(A_1 \times \ldots \times A_n, \set{ \semtop, \sembot }), \\
      &B(f_1: A_1 \to B_1, \ldots, f_n: A_n \to B_n) \coloneqq \parens*
        {
          s \mapsto \parens*
          {
            (b_1, \ldots, b_n) \mapsto
              \begin{rcases}
                \begin{cases}
                  \semtop, &s(a_1, \ldots, a_n) = \semtop \\
                           &\T{and} f_k(a_k) = b_k, \\
                  \sembot, &\T{otherwise}
                \end{cases}
              \end{rcases}
            }
        }
    \end{align*}
    \normalsize

    A natural isomorphism from \( P \) to \( B \) is then
    \begin{equation*}
      \alpha_{(A_1, \ldots, A_n)}(S) \coloneqq \parens*
        {
          (a_1, \ldots, a_n) \mapsto \begin{rcases}
            \begin{cases}
              \semtop, &(a_1, \ldots, a_n) \in S, \\
              \sembot, &\T{otherwise}
            \end{cases}
          \end{rcases}
        }.
    \end{equation*}

    We need to verify that the following diagram commutes:
    \begin{equation*}
      \begin{aligned}
        \includegraphics[page=3]{output/ex__def__natural_transformation}
      \end{aligned}
    \end{equation*}

    Going through the top-right corner gives us
    \begin{align*}
      &\phantom{{}={}}
      \parens[\Big]{ \alpha_{(B_1, \ldots, B_n)} \parens[\big]{ P(f_1, \ldots, f_n)(S) } } (b_1, \ldots, b_n)
      = \\ &=
      \parens[\Big]{ \alpha_{(B_1, \ldots, B_n)} \parens[\big]{ \set[\big]{ \parens[\big]{ f_1(a_1), \ldots, f_n(a_n) } \given (a_1, \ldots, a_n) \in S } } } (b_1, \ldots, b_n)
      = \\ &=
      \begin{cases}
        \semtop, &(a_1, \ldots, a_n) \in S \T{and} f_k(a_k) = b_k \T{for} k = 1, \ldots n, \\
        \sembot, &\T{otherwise},
      \end{cases}
    \end{align*}
    while going through the bottom-left corner gives us
    \begin{align*}
      &\phantom{{}={}}
      \parens[\Big]{ B(f_1, \ldots, f_n) \parens[\big]{ \alpha_{(A_1, \ldots, A_n)}(S) } } (b_1, \ldots, b_n)
      = \\ &=
      \parens*
        {
          (b_1, \ldots, b_n) \mapsto
          \begin{rcases}
            \begin{cases}
              \semtop, &(a_1, \ldots, a_n) \in S \\
                       &\T{and} f_k(a_k) = b_k, \\
              \sembot, &\T{otherwise}
            \end{cases}
          \end{rcases}
        }
      (b_1, \ldots, b_n)
      = \\ &=
      \begin{cases}
        \semtop, &(a_1, \ldots, a_n) \in S \T{and} f_k(a_k) = b_k \T{for} k = 1, \ldots n, \\
        \sembot, &\T{otherwise}.
      \end{cases}
    \end{align*}

    The results coincide, so we conclude that \( \alpha \) is a natural transformation.

    \thmitem{ex:def:natural_transformation/uncurrying_in_set} Another natural isomorphism of sets is \hyperref[con:currying]{function currying} --- the identification of a two-argument function \( s(a, b) \) with its partial application \( a \mapsto (b \mapsto s(a, b)) \).

    More formally, we will construct a natural transformation between the functors
    \begin{align*}
      &L: \cat{Set}^\oppos \times \cat{Set}^\oppos \times \cat{Set} \to \cat{Set}, \\
      &L(A, B, C) \coloneqq \fun(A \times B, C), \\
      &L(f: X \to A, g: Y \to B, h: C \to Z) \coloneqq \parens[\big]{ \underbrace{s}_{\mathclap{A \times B \to C}} \mapsto \underbrace{\parens[\big]{ (x, y) \mapsto h(s(f(x), g(y))) }}_{X \times Y \to Z} }
    \intertext{and}
      &R: \cat{Set}^\oppos \times \cat{Set}^\oppos \times \cat{Set} \to \cat{Set}, \\
      &R(A, B, C) \coloneqq \fun(A, \fun(B, C)), \\
      &R(f: X \to A, g: Y \to A, h: C \to Z) \coloneqq \parens[\big]{ \underbrace{s}_{\mathclap{A \to \fun(B, C)}} \mapsto \underbrace{\parens[\big]{ x \mapsto \parens[\big]{ y \mapsto h\parens[\big]{ \smash{ \overbrace{s(f(x))}^{B \to C}(g(y)) } } } }}_{X \to \fun(Y, Z)} }.
    \end{align*}

    The natural transformation is the obvious one with components
    \begin{equation*}
      \alpha_{(A, B, C)}(s: A \times B \to C) \coloneqq \parens[\big]{ a \mapsto \parens[\big]{ b \mapsto s(a, b) } }.
    \end{equation*}

    Each component is invertible with inverse
    \begin{equation*}
      \alpha_{(A, B, C)}^{-1}(s: A \to \fun(B, C)) \coloneqq \parens[\big]{ (a, b) \mapsto s(a)(b) }.
    \end{equation*}

    It remains to verify the naturality condition
    \begin{equation*}
      \begin{aligned}
        \includegraphics[page=4]{output/ex__def__natural_transformation}
      \end{aligned}
    \end{equation*}

    Indeed, for fixed \( f: X \to A \), \( g: Y \to A \) and \( h: C \to Z \), for any function \( s: A \times B \to C \), going through the top-right corner gives us
    \begin{align*}
      \alpha_{X,Y,Z}\parens[\big]{ L(f, g, h)(s) }
      &=
      \alpha_{X,Y,Z}\parens[\big]{ (x, y) \mapsto h(s(f(x), g(y))) }
      = \\ &=
      \parens[\big]{ x \mapsto y \mapsto h(s(f(x), g(y))) }
    \end{align*}
    while going through the bottom-left corner gives us
    \begin{align*}
      R(f, g, h)\parens[\big]{ \alpha_{A,B,C}(s) }
      &=
      R(f, g, h)\parens[\big]{ a \mapsto b \mapsto s(a, b) }
      = \\ &=
      \parens[\big]{ x \mapsto y \mapsto h(s(f(x), g(y))) }.
    \end{align*}

    This example shows that \( \fun(A, B) \) is an \hyperref[def:exponential_object]{exponential object} in \( \cat{Set} \).

    \thmitem{ex:def:natural_transformation/uncurrying_in_cat} The above uncurrying example extends to the \hyperref[def:category_of_small_categories]{category of small categories} \( \cat{Cat} \), where instead of the set \( \fun(A, B) \) of functions from \( A \) to \( B \) we consider the \hyperref[def:functor]{functor category} \( [\cat{C}, \cat{D}] \) (which we will define soon).

    The functors \( L \) and \( R \) become much more convoluted since they must now act on categories. For example,
    \begin{equation*}
      L(F: \cat{X} \to \cat{A}, G: \cat{Y} \to \cat{B}, H: \cat{C} \to \cat{Z})
    \end{equation*}
    must now itself be a functor, and even
    \begin{equation*}
      L(F, G, H)(S: \cat{A} \times \cat{B} \to \cat{C})
    \end{equation*}
    must be a functor, acting as follows:
    \begin{equation*}
      \begin{aligned}
        &L(F, G, H)(S): \cat{X} \times \cat{Y} \to \cat{Z}, \\
        &L(F, G, H)(S)(X, Y) \coloneqq H(S(F(X), G(Y))), \\
        &L(F, G, H)(S)(f: X \to X', g: Y \to Y') \coloneqq \underbrace{H(S(F(f), G(g)))}_{\mathclap{H(S(F(X), G(Y))) \to H(S(F(X'), G(Y')))}}.
      \end{aligned}
    \end{equation*}

    The functor \( L(F, G, H) \) must also act on natural transformations:
    \begin{equation*}
      L(F, G, H)(\varphi: S \to T) \coloneqq \seq[\big]{ \underbrace{H(\varphi_{(F(X), G(Y))})}_{\mathclap{H(S(F(X), G(Y))) \to H(T(F(X), G(Y)))}} }_{(X, Y) \in \obj(\cat{X} \times \cat{Y}) }.
    \end{equation*}

    The definition of \( R \) is similarly affected. The components of \( \alpha \) must also be functors:
    \begin{equation*}
      \begin{aligned}
        &\alpha_{(\cat{A}, \cat{B}, \cat{C})}: [\cat{A} \times \cat{B}, \cat{C}] \to [\cat{A}, [\cat{B}, \cat{C}]], \\
        &\alpha_{(\cat{A}, \cat{B}, \cat{C})}(S: \cat{A} \times \cat{B} \to \cat{C}) \coloneqq \parens[\big]{ A \mapsto \parens[\big]{ B \mapsto S(A, B) } }, \\
        &\alpha_{(\cat{A}, \cat{B}, \cat{C})}(\varphi: S \to T) \coloneqq \seq{ \underbrace{\seq{ \varphi_{(A, B)} }_{B \in \obj(\cat{B}) }}_{S(A, \anon*) \to T(A, \anon*)} }_{A \in \obj(\cat{A})}
      \end{aligned}
    \end{equation*}

    We can verify the naturality conditions for a natural transformation \( \varphi: S \to T \) in \( [\cat{A} \times \cat{B}, \cat{C}] \): going through the top-right corner gives us
    \begin{align*}
      \alpha_{\cat{X},\cat{Y},\cat{Z}}\parens[\big]{ L(F, G, H)(\varphi) }
      &=
      \alpha_{\cat{X},\cat{Y},\cat{Z}}\parens[\bigg]{ \seq[\big]{ H(\varphi_{(F(X), G(Y))}) }_{(X, Y) \in \obj(\cat{X} \times \cat{Y})} }
      = \\ &=
      \seq{ \seq{ \varphi_{(F(X), G(Y))} }_{Y \in \obj(\cat{Y}) } }_{X \in \obj(\cat{X})}
    \end{align*}
    while going through the bottom-left corner gives us
    \begin{align*}
      R(F, G, H)\parens[\big]{ \alpha_{\cat{A},\cat{B},\cat{C}}(\varphi) }
      &=
      R(F, G, H)\parens[\big]{ \seq{ \seq{ \varphi_{(A, B)} }_{B \in \obj(\cat{B}) } }_{A \in \obj(\cat{A})} }
      = \\ &=
      \seq{ \seq{ \varphi_{(F(X), G(Y))} }_{Y \in \obj(\cat{Y}) } }_{X \in \obj(\cat{X})}.
    \end{align*}

    This example shows that \( [\cat{C}, \cat{D}] \) is an \hyperref[def:exponential_object]{exponential object} in \( \cat{Cat} \).

    \thmitem{ex:def:natural_transformation/directed_multigraphs}\mcite[example 1.3.46]{Perrone2024StartingCategoryTheory} In \cref{def:directed_multigraph}, we have defined a directed multigraph as a set \( V \) of vertices, a set \( A \) of arcs and two functions --- the head \( h: A \to V \) and tail \( t: A \to V \).

    We can regard it as a functor \( G: \cat{I} \to \cat{Set} \), where \( \cat{I} \) is
    \begin{equation*}
      \begin{aligned}
        \includegraphics[page=5]{output/ex__def__natural_transformation}
      \end{aligned}
    \end{equation*}

    Since the index has two objects, a \hyperref[def:natural_transformation]{natural transformation} from the directed multigraph \( G: \cat{I} \to \cat{Set} \) to \( H: \cat{I} \to \cat{Set} \) is then a pair of functions \( f_V: G(V) \to H(V) \) and \( f_A: G(A) \to H(A) \) such that the following \hyperref[def:categorical_diagram]{diagrams commute}:
    \begin{equation*}
      \begin{aligned}
        \includegraphics[page=6]{output/ex__def__natural_transformation}
        \qquad\qquad
        \includegraphics[page=7]{output/ex__def__natural_transformation}
      \end{aligned}
    \end{equation*}

    These are precisely the conditions from \cref{def:directed_multigraph/morphism} characterizing \( (f_V, f_A) \) as a homomorphism of directed multigraphs.
  \end{thmenum}
\end{example}

\begin{definition}\label{def:global_element}\mcite[36]{Awodey2010CategoryTheory}
  In a \hyperref[def:category]{category} with a \hyperref[def:universal_objects/terminal]{terminal object} \( 1 \), we call morphisms \emph{from} \( 1 \) \term{global elements}.
\end{definition}
\begin{comments}
  \item As shown in \cref{ex:def:natural_transformation/global_elements}, global elements generalize elements in concrete categories.
\end{comments}

\begin{proposition}\label{thm:small_natural_transformations}
  A \hyperref[def:natural_transformation]{natural transformation} \( \alpha: F \Rightarrow G \) is a \hyperref[def:small_set]{small set} if and only if each of its components is a small set.
\end{proposition}
\begin{comments}
  \item In case the components are function, this holds for both interpretation of function size from \cref{rem:function_size}.

  Similarly, if the components are functors, this holds for both interpretations of functor size from \cref{rem:functor_size}.
\end{comments}
\begin{proof}
  Follows from \cref{thm:small_set_closure/indexed_family}.
\end{proof}

\paragraph{Functor categories}

\begin{definition}\label{def:functor_category}\mcite[def. 1.4.18]{Perrone2024StartingCategoryTheory}
  For fixed \hyperref[def:category]{categories} \( \cat{C} \) and \( \cat{D} \), we define the \term{functor category} \( [\cat{C}, \cat{D}] \) as follows:
  \begin{thmenum}
    \thmitem{def:functor_category/objects} The \hyperref[def:category/objects]{objects} are, unsurprisingly, the \hyperref[def:functor]{functors} from \( \cat{C} \) to \( \cat{D} \).

    \thmitem{def:functor_category/morphisms} The \hyperref[def:category/morphisms]{morphisms} between two functors are the \hyperref[def:natural_transformation]{natural transformations} between them.

    \thmitem{def:functor_category/composition} The \hyperref[def:category/composition]{composition} of the transformations \( {\alpha: F \Rightarrow G} \) and \( \beta: G \Rightarrow H \) is their componentwise composition, i.e. the transformation \( \gamma: F \Rightarrow H \) with components \( \gamma_A \coloneqq \beta_A \bincirc \alpha_A \).

    We also call this \term{vertical composition} because of how it is visualized in \cref{fig:def:functor_category}.

    \begin{figure}
      \centering
      \includegraphics[page=1]{output/def__functor_category}
      \caption{The \hyperref[def:functor_category/composition]{vertical composition} of \hyperref[def:natural_transformation]{natural transformations}.}\label{fig:def:functor_category}
    \end{figure}

    \thmitem{def:functor_category/identity} The \hyperref[def:category/identity]{identity} on a functor is the natural automorphism whose components are identity morphisms.
  \end{thmenum}
\end{definition}
\begin{defproof}
  \SubProof{Proof that composition is well-defined} The componentwise composition \( \alpha: F \Rightarrow G \) and \( \beta: G \Rightarrow H \) is indeed a natural transformation from \( F \) to \( H \) because
  \begin{equation*}
    \begin{aligned}
      \includegraphics[page=2]{output/def__functor_category}
    \end{aligned}
  \end{equation*}

  \SubProof{Proof that \( [\cat{C}, \cat{D}] \) is a category} The properties \ref{def:category/C1} and \ref{def:category/C2} hold for natural transformations in \( [\cat{C}, \cat{D}] \) since they hold for their components in \( \cat{D} \).
\end{defproof}

\begin{example}\label{ex:def:functor_category}
  We list some examples of \hyperref[def:functor_category]{functor categories}:
  \begin{thmenum}
    \thmitem{ex:def:functor_category/empty} The simplest example is the category \( [\cat{C}, \cat{0}] \) of functors into the \hyperref[def:empty_category]{empty category}. It is itself empty for every \( \cat{C} \).

    This in particular provides a counterexample to the converse of \cref{thm:functor_category_size} --- \( [\cat{C}, \cat{0}] \) is a \hyperref[def:small_category]{small category} even if \( \cat{C} \) is large.

    \thmitem{ex:def:functor_category/initial} For any \( \cat{C} \), the functor category \( [\cat{0}, \cat{C}] \) consists of the empty functor and its identity transformation.

    \thmitem{ex:def:functor_category/unit} Similarly, \( [\cat{1}, \cat{C}] \) consists of the \hyperref[def:global_element]{global elements} in \( \cat{C} \), i.e. the one-object subcategories (with the corresponding identity loop).

    More generally, for any category \( \cat{D} \), \( [\cat{D}, \cat{C}] \) contains the functors that \hyperref[def:factors_through]{factor through} \( \cat{1} \) --- namely, for any object \( A \) in \( \cat{C} \), the functor
    \begin{equation*}
      \begin{aligned}
        &F_A: \cat{D} \to \cat{C}, \\
        &F_A(X) \coloneqq A, \\
        &F_A(f: X \to Y) \coloneqq \id_A.
      \end{aligned}
    \end{equation*}
  \end{thmenum}
\end{example}

\begin{proposition}\label{thm:functor_category_size}
  If both \( \cat{C} \) and \( \cat{D} \) are \hyperref[def:small_category]{small categories} (i.e. \hyperref[def:small_set]{small sets}), so is their \hyperref[def:functor_category]{functor category} \( [\cat{C}, \cat{D}] \).
\end{proposition}
\begin{comments}
  \item \Cref{ex:def:functor_category/empty} provides a simple example to the converse.
\end{comments}
\begin{proof}
  We will use the characterization from \cref{thm:small_set_closure/category} of a small category as a category with a small set of morphisms.

  A natural transformation in \( [\cat{C}, \cat{D}] \) is an indexed family of morphisms in \( \cat{D} \). It is thus an element of the Cartesian product
  \begin{equation*}
    M \coloneqq \prod_{\mathclap{A \in \obj(\cat{C})}} \mor(\cat{D}).
  \end{equation*}

  \Cref{thm:small_set_closure/cartesian_product} implies that \( M \) is a small set if and only if both \( \obj(\cat{C}) \) and \( \mor(\cat{D}) \) are small sets.

  Both latter conditions holds, so \( M \) is a small set, and hence also the set of morphisms in \( [\cat{C}, \cat{D}] \) as a subset.
\end{proof}

\begin{proposition}\label{thm:natural_transformation_cancellation}
  If all components of the \hyperref[def:natural_transformation]{natural transformation} \( \alpha: F \Rightarrow G \) between the pair \( F, G: \cat{C} \to \cat{D} \) are \hyperref[def:morphism_cancellation/left]{monomorphisms} (resp. epimorphisms) in \( \cat{D} \), then \( \alpha \) is a monomorphism (resp. epimorphism) in the \hyperref[def:functor_category]{functor category} \( [\cat{C}, \cat{D}] \).
\end{proposition}
\begin{comments}
  \item Compare this to the analogous invertibility statement \cref{thm:natural_transformation_one_sided_invertibility}.
\end{comments}
\begin{proof}
  Fix a natural transformation \( \alpha: F \Rightarrow G \) between \( F, G: \cat{C} \to \cat{D} \) and suppose that \( \alpha_A \) is a monomorphism in \( \cat{D} \) for every object \( A \) in \( \cat{C} \).

  Then, for any pair of natural transformations \( \beta, \gamma: F \Rightarrow G \) such that
  \begin{equation*}
    \alpha \bincirc \beta = \alpha \bincirc \gamma,
  \end{equation*}
  we have, for any object \( A \) in \( \cat{C} \),
  \begin{equation*}
    \alpha_A \bincirc \beta_A = \alpha_A \bincirc \gamma_A.
  \end{equation*}

  Since \( \alpha_A \) is assumed to be left cancellable, this implies that \( \beta_A = \gamma_A \). Therefore, all components of \( \beta \) and \( \gamma \) coincide, and thus \( \beta = \gamma \).

  We have shown that \( \alpha \) is left cancellable. The case of right cancellation is dual.
\end{proof}

\begin{proposition}\label{thm:natural_transformation_one_sided_invertibility}
  If the \hyperref[def:natural_transformation]{natural transformation} \( \alpha: F \Rightarrow G \) between the pair \( F, G: \cat{C} \to \cat{D} \) is left (resp. right) \hyperref[def:morphism_invertibility]{invertible} in the \hyperref[def:functor_category]{functor category} \( [\cat{C}, \cat{D}] \), then its components are left (resp. right) invertible in \( \cat{D} \).
\end{proposition}
\begin{comments}
  \item Compare this to the analogous cancellation statement \cref{thm:natural_transformation_cancellation}.
  \item The converse is more subtle --- see the proof of \cref{thm:natural_isomorphism}.
\end{comments}
\begin{proof}
  Let \( \beta: G \Rightarrow F \) be a left inverse of \( \alpha: F \Rightarrow G \) Then, for each object \( A \) in \( \cat{C} \), we have
  \begin{equation*}
    \beta_A \bincirc \alpha_A = \id_A,
  \end{equation*}
  i.e. \( \beta_A \) is a left inverse of \( \alpha_A \) in \( \cat{D} \).

  Similarly, if \( \beta \) is a right inverse of \( \alpha \), then \( \beta_A \) is a right inverse of \( \alpha_A \).
\end{proof}

\begin{proposition}\label{thm:natural_isomorphism}
  The \hyperref[def:natural_transformation]{natural transformation} \( \alpha: F \Rightarrow G \) between the pair \( F, G: \cat{C} \to \cat{D} \) is an \hyperref[def:morphism_invertibility/isomorphism]{isomorphism} in the \hyperref[def:functor_category]{functor category} \( [\cat{C}, \cat{D}] \) if and only if all its components are isomorphisms in \( \cat{D} \).

  Furthermore, the family \( \seq{ \alpha_A^{-1} }_{A \in \obj(\cat{C})} \) of inverse components of \( \alpha \) is a natural transformation that acts as the inverse of \( \alpha \) in \( [\cat{C}, \cat{D}] \).
\end{proposition}
\begin{comments}
  \item This proposition justifies the definition for natural isomorphisms in \cref{def:natural_isomorphism}.
\end{comments}
\begin{proof}
  \SufficiencySubProof If \( \beta \) has a two-sided inverse of \( \beta \) in \( [\cat{C}, \cat{D}] \), then \( \beta_A \) is, by \cref{thm:natural_transformation_one_sided_invertibility}, both a left and right inverse of \( \alpha_A \) in \( \cat{D} \). \Cref{thm:def:morphism_invertibility/left_and_right} implies that \( \beta_A \) is the unique two-sided inverse of \( \alpha_A \).

  \NecessitySubProof Suppose that, for every object \( A \) in \( \cat{C} \), \( \alpha_A \) has a two-sided inverse \( \beta_A \). We can pack these inverses into a single family \( \beta \).

  It is immediate that \( \beta \) is the two-sided inverse of \( \alpha \), \emph{as long as} it is a natural transformation. We must verify that the following diagram commutes:
  \begin{equation}\label{eq:thm:natural_isomorphism/proof/necessity_diagram}
    \begin{aligned}
      \includegraphics[page=1]{output/thm__natural_isomorphism}
    \end{aligned}
  \end{equation}

  This is part of the following larger diagram, the rest of which are our assumptions:
  \begin{equation*}
    \begin{aligned}
      \includegraphics[page=2]{output/thm__natural_isomorphism}
    \end{aligned}
  \end{equation*}

  Naturality of \( \alpha \) gives us the equality (from the diagram above):
  \begin{equation*}
    \alpha_B \bincirc F(f) = G(f) \bincirc \alpha_A.
  \end{equation*}

  Since \( \beta_B \bincirc \alpha_B = \id_{F(f)} \), composing both sides with \( \beta_B \) gives us
  \begin{equation*}
    F(f) = \beta_B \bincirc G(f) \bincirc \alpha_A.
  \end{equation*}

  Since \( \alpha_A \bincirc \beta_A = \id_{G(f)} \), precomposing both sides with \( \beta_A \) gives us
  \begin{equation*}
    F(f) \bincirc \beta_A = \beta_B \bincirc G(f).
  \end{equation*}

  This is precisely the commutativity condition for \eqref{eq:thm:natural_isomorphism/proof/necessity_diagram}.
\end{proof}

\paragraph{Opposite natural transformations}

\begin{remark}\label{rem:opposite_natural_transformation}
  Fix a pair of parallel functors \( F, G: \cat{C} \to \cat{D} \) and consider some natural transformation \( \alpha: F \to G \). By definition, it must satisfy the naturality condition
  \begin{equation*}
    \begin{aligned}
      \includegraphics[page=1]{output/rem__opposite_natural_transformation}
    \end{aligned}
  \end{equation*}

  For the \hyperref[def:opposite_functor]{opposite functor} \( F^\oppos \), given \( f: A \to B \), we have defined \( F^\oppos(f) \) as \( F(f)^* \), thus,
  \begin{equation*}
    F^\oppos(f): F^\oppos(B) \to F^\oppos(A).
  \end{equation*}

  Moreover, since \( \alpha_A: F(A) \to G(A) \), we have
  \begin{equation*}
    \alpha_A^\oppos: G^\oppos(A) \to F^\oppos(A).
  \end{equation*}

  Since the above diagram commutes, so does
  \begin{equation*}
    \begin{aligned}
      \includegraphics[page=2]{output/rem__opposite_natural_transformation}
    \end{aligned}
  \end{equation*}

  This demonstrates the naturality of \( \alpha^\oppos: G^\oppos \to F^\oppos \), the family of morphisms defined by inverting the components of \( \alpha \).
\end{remark}

\begin{definition}\label{def:opposite_natural_transformation}\mcite[144]{ЦаленкоШульгейфер1974ОсновыТеорииКатегорий}
  We define the \term[ru=двойственное (естественное преобразование)]{opposite} of a \hyperref[def:natural_transformation]{natural transformation} \( \alpha: F \to G \) by taking the opposites of all its components, thus obtaining \( \alpha^\oppos: G^\oppos \to F^\oppos \).
\end{definition}
\begin{defproof}
  We show in \cref{rem:opposite_natural_transformation} that the naturality diagram for \( \alpha^\oppos \) commutes.
\end{defproof}

\begin{proposition}\label{thm:opposite_natural_transformation}
  Fix a pair \( F, G: \cat{C} \to \cat{D} \) of parallel functors. Fix also a family
  \begin{equation*}
    \alpha = \seq{ \alpha_A }_{A \in \obj(\cat{C})}
  \end{equation*}
  of morphisms in \( \cat{D} \) such that \( \alpha_A: F(A) \to G(A) \).

  Then \( \alpha \) is a \hyperref[def:natural_transformation]{natural transformation} from \( F \) to \( G \) if and only if its \hyperref[def:opposite_natural_transformation]{opposite} \( \alpha^\oppos \) is natural transformation from \( G^\oppos \) to \( F^\oppos \).
\end{proposition}
\begin{proof}
  Follows from \cref{rem:opposite_natural_transformation}.
\end{proof}
